\section*{Discussion}
In this study, we perform a large-scale empirical measurement of the effect of adopting AI in science on both individual scientists and scientific communities. We identify three waves of AI adoptions in science, corresponding with the dominance of machine learning, deep learning, and large language models. Each wave is marked with an accelerated AI adoption rate in research papers and authors.
We find that individual scientists are increasingly rewarded with expanded academic impact and accelerated career development for incorporating AI assistance in research across these waves and in all natural science research fields we studied.
On average, the use of AI helps individual scientists publish 3.02 times more papers, receive 4.84 times more citations, and become team leaders 1.37 years earlier. This substantial academic benefit may be a driving force behind the accelerated rate of AI adoption.
However, we also find unintended consequences from the increased prevalence of AI-augmented research.
In all fields, AI-augmented research focuses on a narrower scope of scientific topics and reduces the scientific engagement of follow-on research, leading to more overlapping research works that slows the expansion of knowledge.
Further, with a greater concentration of collective attention to the same AI papers, the adoption of AI appears to induce authors to engage in collective hill-climbing\footnote{The metaphor of “collective hill-climbing” describes a situation where researchers act like a group of climbers all scaling the same popular mountain from the same route. Because both the “path” of a known approach and the “peak” of an anticipated solution are constrained, this collective rush leads to “crowding” and may discourage the search for other, potentially higher “mountains” representing new questions and answers.}, catalyzing solutions to known problems rather than creating new ones.

These findings raise critical questions for science policy. What are the topics most likely left behind from AI-augmented research across fields? Those with less available data include critical scientific questions regarding the origins of natural phenomena, where data are necessarily reduced. Accelerating scientific activity “under the lamp post” of highly visible, data-rich phenomena moves science away from many foundational questions and towards operational ones. By driving attention toward the most popular new developments, AI appears to drive problem solution over generation. These issues become particularly concerning in the face of calls to further increase support for AI-augmented science~\supercite{borger2023artificial,lawrence2024accelerating}, coupled with the personal scientific incentives we demonstrate. This could shift collective attention away from new and original questions that lack the data required for AI to demonstrate benefit.
It is true that more overlapping attention and a contracted focus may benefit scientific replication and extension, accelerating the emergence of solid and practical solutions to specific questions.
Insofar as scientific discovery represents a vast and complex landscape, however, concentrating attention on the same developments may increase the likelihood that science becomes fixed on local maxima of scientific explanation and prediction rather than searching in a more broad, decoupled, and diverse way.

While our analysis provides new insight into AI's impact on science, clear limitations remain. Our identification approach, though validated by experts, misses subtle and unmentioned forms of AI use, and our focus on natural sciences excludes important domains where AI adoption patterns may differ. Moreover, despite consistently suggestive evidence, we cannot fully identify the causal linkage between AI adoption and scientific impact. Nevertheless, our findings demonstrate that currently attributed uses of AI in science primarily augment cognitive tasks through data processing and pattern recognition. Looking forward, these findings illuminate a critical and expansive pathway for AI development in science. To preserve collective exploration in an era of AI use, we will need to reimagine AI systems that expand not only cognitive capacity but also sensory and experimental capacity~\supercite{king2009automation,burger2020mobile}, enabling scientists to search, select, and gather new types of data from previously inaccessible domains rather than merely optimizing analysis of standing data. The history of major discoveries has been most consistently linked with new views on nature~\supercite{krauss2024debunking}.
Expanding the scope of AI's deployment in science will be required for sustained scientific research and to stimulate new fields rather than merely automate existing ones.